% Ganti setiap \titik dengan poin rangkuman praktikum.
% Word Heading2b: Arial bold italic, isi Garamond dengan bullet kotak.
\vspace*{-8.8mm}
{\centering\bfseries\fontsize{15}{18}\selectfont\judul\par}\vspace{1.4mm}
\begin{rangkuman}
  \linespread{1.0}\selectfont
  {\arialfont\bfseries\itshape\fontsize{12}{14}\selectfont Rangkuman praktikum}\par\vspace{4mm}
  {\arialfont\bfseries\itshape\fontsize{10}{12}\selectfont
  Poin-poin utama yang anda ketahui selama praktikum adalah:\par}\vspace{-1.5mm}
  \begin{itemize}[label=\raisebox{0.45ex}{\fontsize{3.2}{3.2}\selectfont$\blacksquare$},leftmargin=10.9mm,labelsep=3.85mm,labelwidth=0pt,
                  align=left,itemsep=0pt,topsep=0pt,parsep=0pt,
                  before=\normalfont\fontsize{9.6}{17.4}\selectfont]
    \item \titik
    \item \titik
    \item \titik
    \item \titik
    \item \titik
  \end{itemize}
\end{rangkuman}
\section{Dokumentasi Praktikum}

\subsection{Jelaskan Langkah-langkah yang dilakukan selama praktikum}

\begin{itemize}
  \item Langkah pertama adalah (silahkan jelaskan langkah-langkah
      selama praktikum)

      Lampirkan gambar yang dilakukan selama langkah pertama.

      \begin{figure}[H]
        \centering
        \gambarpengganti{figure/contoh-gambar.png}
        \caption{Ilustrasi dokumentasi praktikum}
      \end{figure}
  \newpage % sama dengan Word: langkah kedua mulai di halaman baru (hapus jika tidak perlu)
  \item Langkah kedua adalah (silahkan jelaskan langkah-langkah
      selama praktikum)

      Jika ada tabel yang perlu dilampirkan berikut format table nya.

      \begin{table}[H]
        \centering
        \caption{Contoh penggunaan table}
        \linespread{1}\fontsize{8}{10}\selectfont
        \rowcolors{2}{tabband}{white}
        \setlength{\tabcolsep}{3pt}
        \renewcommand{\arraystretch}{0.9}
        \arrayrulecolor{tabline}
        \begin{tabular}{!{\color{tabline}\vrule}>{\centering\arraybackslash}p{9mm}!{\color{tabline}\vrule}p{30mm}!{\color{tabline}\vrule}p{30mm}!{\color{tabline}\vrule}p{30mm}!{\color{tabline}\vrule}}
          \hline
          \rowcolor{tabhead}
          \thead{No} & \thead{Judul kolom 1} & \thead{Judul Kolom 2} & \thead{Judul Kolom 3} \\
          \hline
          & & & \\ \hline
          & & & \\ \hline
          & & & \\ \hline
          & & & \\ \hline
          & & & \\ \hline
          & & & \\ \hline
          & & & \\ \hline
          & & & \\ \hline
        \end{tabular}
      \end{table}
  \item Langkah ketiga adalah (silahkan jelaskan langkah-langkah
      selama praktikum). Jika ada penggunaan list atau item di dalam
      penjelasan anda, silahkan gunakan format bullet and numbering berikut:
      \begin{itemize}
        \item List 1
        \item List 2
        \item List 3
      \end{itemize}
\end{itemize}
\par
\subsubsection{Alat dan Bahan}

Jelaskan alat dan bahan yang digunakan dalam praktikum. Deskripsikan alat atau
device yang anda gunakan. Jika menggunakan komputer Windows, cukup dengan
memberikan screenshoot dxdiag.
\par
\par
\subsubsection{Tugas}

Silahkan jelaskan jawaban dari tugas yang diberikan. Jika tidak ada tugas, maka
setiap praktikum anda wajib melampirkan dokumen laporan praktikum ini.
\par
\nocite{forouzan2007}




